% Định dạng dùng chung cho báo cáo. Biên dịch bằng XeLaTeX.
\usepackage[margin=2.54cm,headheight=15pt]{geometry}
\usepackage{fontspec}
\IfFontExistsTF{Times New Roman}
  {\setmainfont{Times New Roman}}
  {\setmainfont{TeX Gyre Termes}}
\IfFontExistsTF{Consolas}
  {\setmonofont{Consolas}[Scale=MatchLowercase]}
  {\setmonofont{Latin Modern Mono}[Scale=MatchLowercase]}
\usepackage{setspace}
\usepackage{graphicx}
\usepackage{pdflscape}
\usepackage{booktabs,tabularx,longtable}
\usepackage{xcolor}
\usepackage{fvextra,needspace}

\usepackage[hidelinks,unicode]{hyperref}

\onehalfspacing
\setlength{\parindent}{1.27cm}
\setlength{\parskip}{0pt}
\emergencystretch=1.5em
\widowpenalty=10000
\clubpenalty=10000
\displaywidowpenalty=10000
\raggedbottom

\renewcommand{\contentsname}{MỤC LỤC}
\renewcommand{\figurename}{Hình}
\renewcommand{\tablename}{Bảng}
\setcounter{tocdepth}{2}
\setcounter{secnumdepth}{3}

\makeatletter
\newcommand{\ps@report}{%
  \let\@oddfoot\@empty
  \let\@evenfoot\@empty
  \def\@oddhead{\hfill\normalfont\thepage}%
  \let\@evenhead\@oddhead
}
\pagestyle{report}

% Tiêu đề chương có tiền tố trong thân bài, mục lục vẫn dùng số 1, 2, ...
\newcommand{\ReportSectionPrefix}{Chương }
\renewcommand{\@seccntformat}[1]{%
  \ifcsname report@prefix@#1\endcsname\csname report@prefix@#1\endcsname\fi
  \csname the#1\endcsname
  \ifcsname report@prefix@#1\endcsname.\enspace\else\quad\fi
}
\let\report@prefix@section\ReportSectionPrefix

% Khoảng cách trước đề mục dương: thụt đầu dòng cả đoạn sau đề mục.
\renewcommand{\section}{\@startsection{section}{1}{\z@}%
  {14pt plus 3pt minus 2pt}{5pt plus 1pt}%
  {\normalfont\fontsize{16}{20}\selectfont\bfseries\raggedright}}
\renewcommand{\subsection}{\@startsection{subsection}{2}{\z@}%
  {10pt plus 2pt minus 1pt}{3pt plus 1pt}%
  {\normalfont\fontsize{14}{17}\selectfont\bfseries\raggedright}}
\renewcommand{\subsubsection}{\@startsection{subsubsection}{3}{\z@}%
  {8pt plus 2pt minus 1pt}{2pt plus 1pt}%
  {\normalfont\normalsize\bfseries\raggedright}}

\renewcommand{\tableofcontents}{%
  \section*{\centering\contentsname}%
  \begingroup
    \singlespacing\small
    \renewcommand{\@pnumwidth}{2.5em}%
    \renewcommand{\@tocrmarg}{3.5em}%
    \@starttoc{toc}%
  \endgroup
}
\renewcommand{\l@section}[2]{%
  \addpenalty\@secpenalty
  \addvspace{3pt plus 1pt}%
  \@dottedtocline{1}{0em}{1.85em}{\bfseries #1}{\bfseries #2}%
}
\renewcommand{\l@subsection}{\@dottedtocline{2}{1.85em}{2.7em}}
\newcommand{\l@reportappendix}[2]{%
  \addpenalty\@secpenalty
  \addvspace{3pt plus 1pt}%
  \@dottedtocline{1}{0em}{5.6em}{\bfseries #1}{\bfseries #2}%
}
\newcommand{\toclevel@reportappendix}{1}

\newcommand{\ReportListOfFigures}{%
  \ReportSection{Danh mục hình}%
  {\small\singlespacing\@starttoc{lof}}%
}
\newcommand{\ReportListOfTables}{%
  \ReportSection{Danh mục bảng}%
  {\small\singlespacing\@starttoc{lot}}%
}

% Chú thích theo APA cho figure, table và longtable.
\newcommand{\ReportCaptionText}[2]{%
  \raggedright\setlength{\parindent}{0pt}%
  #1\textit{#2}\par
}
\renewcommand{\@makecaption}[2]{%
  \vskip\abovecaptionskip
  \noindent\parbox[t]{\linewidth}{%
    \small\singlespacing\ReportCaptionText{\textbf{#1}\par}{#2}}\par
  \vskip\belowcaptionskip
}
\renewcommand{\LT@makecaption}[3]{%
  \LT@mcol\LT@cols l{\hbox to\z@{%
    \parbox[t]{\LTcapwidth}{\small\singlespacing
      \ReportCaptionText{#1{\textbf{#2}\par}}{#3}\medskip}%
    \hss}}%
}
% Giữ hình và thuyết minh cùng thứ tự đọc khi float có thể trôi quá xa.
\newenvironment{ReportFigureHere}{%
  \par\medskip\noindent\begin{minipage}{\linewidth}\def\@captype{figure}%
}{%
  \end{minipage}\par\medskip
}
\newenvironment{ReportTableHere}{%
  \par\medskip\noindent\begin{minipage}{\linewidth}\def\@captype{table}%
}{%
  \end{minipage}\par\medskip
}
\makeatother

\setlength{\abovecaptionskip}{0pt}
\setlength{\belowcaptionskip}{7pt}
\setlength{\LTcapwidth}{\textwidth}
\setlength{\LTleft}{0pt}
\setlength{\LTright}{0pt plus 1fill}
\newcolumntype{L}{>{\raggedright\arraybackslash}X}
\newcolumntype{P}[1]{>{\raggedright\arraybackslash}p{#1}}
\newcommand{\ReportNote}[1]{%
  \par\smallskip{\small\singlespacing\raggedright
    \noindent\textit{Ghi chú.} #1\par}%
}
\newcommand{\ReportImage}[2][]{%
  {\centering\includegraphics[#1,width=\linewidth,height=.68\textheight,keepaspectratio]{#2}\par}%
}
\newcommand{\ReportUmlGraphic}[2]{%
  {\centering
   \edef\ReportUmlOptions{width=\the\linewidth,height=\the\dimexpr#1\textheight\relax,
     keepaspectratio,clip,trim=\csname ReportUmlTrim-#2\endcsname}%
   \expandafter\includegraphics\expandafter[\ReportUmlOptions]{report/figures/#2.pdf}\par}%
}
\newcommand{\ReportUmlTrim}[2]{\expandafter\def\csname ReportUmlTrim-#1\endcsname{#2}}
\newcommand{\ReportUml}[4][.90]{%
  \begin{figure}[!htbp]
    \caption{#4}\label{#3}%
    \ReportUmlGraphic{#1}{#2}%
  \end{figure}%
}
\newcommand{\ReportWideUml}[3]{%
  \begin{landscape}
    \thispagestyle{empty}%
    \begin{ReportFigureHere}
      \vbox to 0pt{\vss\hbox to \linewidth{\hfill\thepage}\kern 12pt}%
      \nointerlineskip
      \caption{#3}\label{#2}%
      \ReportUmlGraphic{.90}{#1}%
    \end{ReportFigureHere}
  \end{landscape}%
}
% Giữ chuỗi Unicode đầu vào khi sao chép từ PDF, kể cả dấu - của Consolas.
\newcommand{\ReportTechnicalText}{\XeTeXgenerateactualtext=1\relax}
\newcommand{\code}[1]{\texttt{\ReportTechnicalText\detokenize{#1}}}

% XeLaTeX đọc trực tiếp Unicode trong khối verbatim.
\newcommand{\ReportMonoSize}{\fontsize{10}{12}\selectfont}
\newcommand{\ReportCodeSize}{\fontsize{9.5}{11.5}\selectfont}
\definecolor{ReportCodeFocus}{gray}{0.94}
\fvset{
  fontfamily=tt,
  fontsize=\ReportMonoSize,
  formatcom=\ReportTechnicalText,
  baselinestretch=1,
  tabsize=4,
  obeytabs=true,
  breaklines=true,
  breakanywhere=false,
  breakautoindent=true,
  breakindent=1.5em,
  breaksymbolleft={},
  highlightcolor=ReportCodeFocus,
  listparameters={\setlength{\topsep}{4pt}\setlength{\partopsep}{0pt}}
}
\renewcommand{\theFancyVerbLine}{\scriptsize\arabic{FancyVerbLine}}
% Trong đoạn trích có commandchars=§«», giữ literal liền nhau và cho phép ngắt trước nó.
\newcommand{\ReportLiteral}[1]{\allowbreak\mbox{#1}}

\newcounter{reportcode}
\newcommand{\ReportCodeHeading}[2]{%
  \par\Needspace{5\baselineskip}%
  \refstepcounter{reportcode}\label{#2}%
  \smallskip
  {\small\singlespacing\raggedright\noindent
    \textbf{Mã nguồn \thereportcode}\par\nopagebreak
    \noindent\textit{#1}\par}%
  \nopagebreak
}
\newenvironment{ReportCode}[3][]{%
  \VerbatimEnvironment
  \ReportCodeHeading{#2}{#3}%
  \begin{Verbatim}[fontsize=\ReportCodeSize,numbers=left,xleftmargin=1.7em,numbersep=6pt,#1]%
}{\end{Verbatim}}
\newcommand{\ReportCodeFile}[4][]{%
  \ReportCodeHeading{#2}{#3}%
  \VerbatimInput[fontsize=\ReportCodeSize,numbers=left,xleftmargin=1.7em,numbersep=6pt,#1]{#4}%
}
\DefineVerbatimEnvironment{ReportTerminal}{Verbatim}{numbers=none}
\DefineVerbatimEnvironment{ReportJson}{Verbatim}{numbers=none}
\DefineVerbatimEnvironment{ReportOutput}{Verbatim}{numbers=none}
\newcommand{\ReportCodeRef}[1]{Mã nguồn~\ref{#1}}
\newcommand{\ReportCodeSource}[1]{%
  \par\smallskip{\small\singlespacing\noindent\textit{Nguồn mã.} #1\par}%
}

\newcommand{\ReportCover}[1]{%
  \hypersetup{pageanchor=false}%
  \begingroup
    \pagestyle{empty}\thispagestyle{empty}%
    \input{#1}\clearpage
  \endgroup
  \hypersetup{pageanchor=true}%
}
\newcommand{\ReportChapter}[1]{\clearpage\section{#1}}

% Đề mục không đánh số vẫn xuất hiện trong mục lục.
\newcommand{\ReportSection}[1]{%
  \section*{#1}%
  \addcontentsline{toc}{section}{#1}%
}
\newcommand{\ReportReferenceHeading}{%
  \section*{\centering TÀI LIỆU THAM KHẢO}%
  \addcontentsline{toc}{section}{Tài liệu tham khảo}%
}
% Khai báo mỗi nguồn một lần để trích dẫn và tạo mục tương ứng trong danh mục.
\makeatletter
\newcommand{\ReportSourceList}{}
\newcommand{\ReportDeclareSource}[4]{%
  \@namedef{report@cite@#1}{\hyperlink{source:#1}{#2, #3}}%
  \g@addto@macro\ReportSourceList{\ReportSourceEntry{#1}{#4}}%
}
\newcommand{\ReportCite}[1]{(\@nameuse{report@cite@#1})}
\makeatother
\newcommand{\ReportSourceEntry}[2]{%
  \par\noindent\hangindent=1.27cm\hangafter=1%
  \hypertarget{source:#1}{}#2\par%
}
% Giữ mỗi URL đủ ngắn trên một dòng và dùng văn bản gốc khi sao chép PDF.
\newcommand{\ReportReferenceURL}[1]{\href{#1}{\mbox{\code{#1}}}}
\newcommand{\ReportBibliography}{%
  \clearpage
  \ReportReferenceHeading
  \ReportSourceList
}

% Bắt đầu phụ lục mới, giữ A.1, A.1.1 cho các đề mục con.
\newcommand{\ReportAppendix}[1]{%
  \clearpage
  \refstepcounter{section}%
  \section*{Phụ lục \thesection. #1}%
  \addcontentsline{toc}{reportappendix}{\protect\numberline{Phụ lục \thesection}#1}%
}
